% SPDX-License-Identifier: CC-BY-4.0
\section{Introduction}
\label{sec:intro}

Behavioral host intrusion detection systems watch a process's
syscalls, file access, network endpoints, and child spawns, fit a
per-process baseline from those observations, and score live behavior
as drift against the baseline. The lineage starts with Forrest's
``sense of self''~\cite{forrest1996sense} and runs through n-gram
sequence models, hidden-Markov detectors, and modern provenance-graph
systems~\cite{goyal2023mimicry}. The features changed; the assumption
that the learning window is adversary-free did not.

That assumption fails in practice. Mandiant's M-Trends 2025 report
~\cite{mandiant2025mtrends} puts the global median attacker dwell time
at 11 days for 2024 incidents, with a longer tail for sophisticated
threat actors. Any detector that turns on while the adversary is
already on the host fits a baseline shaped by the adversary. After
that, the adversary's behavior \emph{is} the baseline's notion of
normal, and the detector cannot tell planted behavior apart from
genuine drift.

We call this \emph{baseline poisoning}. The lever is not exotic. A
single elevated shell session can shape the baseline of every
behavioral feature the detector tracks for that process, and a
carefully-paced injection slips past any change-point detector
watching the live stream. The question is what property an estimator
needs to stay useful when the learning window is partially
adversarial, and what an operator pays for that property in
clean-condition false-positive rate.

This paper makes four contributions.

\begin{enumerate}
\item We \textbf{formalize} the baseline-poisoning threat model
  (\cref{sec:threat-model}). The adversary has a budget
  $\budget \in [0, 1]$ that bounds the fraction of the learning
  window they control. We consider observability variants up to a
  white-box adversary that knows the defender's hyperparameters.
\item We \textbf{demonstrate} that the naive mean-and-stddev
  estimator widely used in prior work loses roughly a factor of
  3.6 in target-score discriminability at a 15\% poisoning budget on
  our isotropic synthetic workload, and continues to degrade
  smoothly as the budget grows
  (\cref{sec:eval-naive-collapse}).
\item We \textbf{propose} \sediment{} (\cref{sec:sediment}), a
  per-feature trimmed-moment estimator with a closed-form
  Gaussian-truncation bias correction. Setting $\trim{} \geq 2\budget$
  gives a defensible robustness guarantee against any adversary with
  budget at most \budget{}.
\item We \textbf{evaluate} the design space across nine experiments
  (\cref{sec:evaluation}), including a four-way ablation against
  alternative robustness primitives and three poisoning attacks:
  uniform shuffled, contiguous burst, and a grid-search white-box
  adversary that knows \trim{}.
\end{enumerate}

A reference Python implementation accompanies a production Rust
runtime; cross-implementation parity tests on every push keep them
bit-equivalent across both the estimator and the syscall n-gram
extractor. The artifact is open
source.\footnote{\url{https://github.com/ChronosynD/core}}
